\documentclass[11pt]{article}
\usepackage[margin=0.92in]{geometry}
\usepackage{amsmath,amssymb,amsthm}
\usepackage{microtype}
\usepackage{xcolor}
\usepackage{hyperref}
\hypersetup{colorlinks=true,linkcolor=blue!55!black,urlcolor=blue!55!black,citecolor=blue!55!black}
\newtheorem{theorem}{Theorem}
\newtheorem{lemma}[theorem]{Lemma}
\newtheorem{corollary}[theorem]{Corollary}
\newcommand{\LP}{\mathcal{L\!P}}
\title{Parameter monotonicity and a sharp threshold for\\Nguyen's special entire function}
\author{Candidate full-solution packet}
\date{11 August 2026}

\begin{document}
\maketitle

\begin{center}
\fbox{\parbox{0.92\linewidth}{\textbf{Status:} candidate full proof likely valid.  The paper's open parameter-monotonicity conjecture has a positive answer.  Consequently membership is governed by one threshold $a_0$, and an exact rational certificate gives $3.9642<a_0<3.9643$.  The argument uses the paper's Theorem~1, but not its erroneous numerical Theorem~2(ii).}}
\end{center}

\section{Question and result}

For $a>1$ let
\[
 F_a(z)=\sum_{k=0}^{\infty}\frac{z^k}{\prod_{j=1}^k(a^j+1)}.
\]
Nguyen proved that $F_a\in\LP$ if and only if there is an
$x\in(a+1,a^2+1)$ with $F_a(-x)\leq0$~\cite[Theorem~1]{Nguyen}.
The paper then conjectured that membership is monotone in the parameter:
if $F_{a_1}\in\LP$ and $a_2>a_1$, must $F_{a_2}\in\LP$?

\begin{theorem}[Positive answer to the conjecture]\label{thm:main}
If $1<a_1<a_2$ and $F_{a_1}\in\LP$, then $F_{a_2}\in\LP$.
More precisely, set
\[
 q_2(a)=\frac{a^2+1}{a+1},\qquad
 H_a(y)=F_a(-(a+1)y).
\]
For every $a\geq(3+\sqrt{17})/2$ and $0<y\leq q_2(a)$,
\[
 \frac{\partial}{\partial a}H_a(y)<0.
\]
\end{theorem}

\section{Proof Intuition}

The normalized variable $y$ is the key point.  It freezes the left endpoint
of the sign-detecting interval at $1$, while $q_2(a)$ is increasing.  Thus a
single negative witness for $a_1$ remains admissible and becomes strictly
more negative for every larger parameter.  Analytically, differentiation in
$a$ produces an alternating series.  Pairing each even term with its following
odd term leaves a negative pair: the exceptional first pair reduces to one
explicit positive polynomial, while every later pair is dominated by the
exponential growth of $a^{k+1}+1$.  This is the monotonicity mechanism.

\section{Proof of parameter monotonicity}

Write
\[
 H_a(y)=\sum_{k=0}^{\infty}(-1)^k A_k(a)y^k,
 \qquad
 A_k(a)=\frac{(a+1)^k}{\prod_{j=1}^k(a^j+1)}.
\]
Here $A_0=A_1=1$.  For $j\geq2$ define
\[
 \delta_j(a)=\frac{j a^{j-1}}{a^j+1}-\frac1{a+1},
 \qquad
 \lambda_k(a)=\sum_{j=2}^k\delta_j(a).
\]
Since
\[
 j a^{j-1}(a+1)-(a^j+1)
 =(j-1)a^j+j a^{j-1}-1>0,
\]
all $\delta_j$ and $\lambda_k$ are positive.  Logarithmic differentiation
gives $A_k'=-\lambda_kA_k$ for $k\geq2$.  Termwise differentiation is
legitimate by locally uniform absolute convergence, and pairing consecutive
terms gives
\begin{equation}\label{eq:pairs}
 \partial_aH_a(y)=\sum_{\substack{k\geq2\\k\ \mathrm{even}}}
 \left(-\lambda_kA_ky^k+\lambda_{k+1}A_{k+1}y^{k+1}\right).
\end{equation}

We show that every parenthesis is negative.  Because
\[
 \frac{A_{k+1}}{A_k}=\frac{a+1}{a^{k+1}+1}
 \quad\text{and}\quad y\leq q_2(a),
\]
it is enough to prove
\begin{equation}\label{eq:ratio-goal}
 \frac{\lambda_{k+1}}{\lambda_k}
 <\frac{a^{k+1}+1}{a^2+1}.
\end{equation}

The first pair, $k=2$, is the only close one.  Direct simplification shows
that \eqref{eq:ratio-goal} is equivalent to $P(a)>0$, where
\[
 P(a)=a^7-2a^6-4a^5+2a^4-6a^3+a^2-a+1.
\]
For $a=u+7/2$ its coefficients, from constant term upward, are
\[
 \frac{90771}{128},\ \frac{236319}{64},\ \frac{156727}{32},\
 \frac{49107}{16},\ \frac{8521}{8},\ \frac{845}{4},\
 \frac{45}{2},\ 1.
\]
Hence $P(a)>0$ for $a\geq7/2$.

For $k\geq4$, observe that
\[
 \delta_2\geq\frac1{a+1},\qquad
 \delta_{k+1}<\frac{k+1}{a}.
\]
Consequently
\[
 \frac{\lambda_{k+1}}{\lambda_k}
 =1+\frac{\delta_{k+1}}{\lambda_k}
 \leq1+\frac{\delta_{k+1}}{\delta_2}<2k+3.
\]
For $a\geq3$ we also have
\[
 \frac{a^{k+1}+1}{a^2+1}>2k+3\qquad(k\geq4).
\]
Indeed the case $k=4$ reduces to $a^5-11a^2-10>0$, true at
$a=3$ and thereafter by differentiation; and the left side more than
doubles when $k$ is increased by one, which proves the assertion by
induction.  This proves \eqref{eq:ratio-goal} for every pair in
\eqref{eq:pairs}, and therefore $\partial_aH_a(y)<0$ whenever
$a\geq7/2$ and $0<y\leq q_2(a)$.

It remains only to place a possible Laguerre--P\'olya parameter in this
range.  The standard coefficient estimate reproduced as Lemma~1 in
\cite{Nguyen} says that, for this increasing sequence of second quotients,
$F_a\in\LP$ forces $q_2(a)\geq3$.  Thus
\[
 a\geq\frac{3+\sqrt{17}}2>\frac72.
\]
Now suppose $F_{a_1}\in\LP$.  The sign criterion supplies
$y_0\in(1,q_2(a_1))$ with $H_{a_1}(y_0)\leq0$.  Since $q_2$ is strictly
increasing, $y_0<q_2(a)$ throughout $[a_1,a_2]$.  The inequality just proved
therefore gives
\[
 H_{a_2}(y_0)<H_{a_1}(y_0)\leq0.
\]
Applying the sign criterion once more yields $F_{a_2}\in\LP$, proving
Theorem~\ref{thm:main}.

\section{Existence and exact localization of the threshold}

\begin{corollary}\label{cor:threshold}
There is a number $a_0>1$ such that
\[
 F_a\in\LP\quad\Longleftrightarrow\quad a\geq a_0.
\]
Moreover,
\[
 \boxed{\ 3.9642<a_0<3.9643\ }.
\]
\end{corollary}

\begin{proof}
Theorem~\ref{thm:main} makes the membership set upward closed, and it is
nonempty by Hutchinson's $q_n\geq4$ criterion.  It is also closed.  Indeed,
if $a_n$ decreases to its infimum and $F_{a_n}\in\LP$, choose sign witnesses
$y_n\in(1,q_2(a_n))$.  A subsequence converges to
$y\in[1,q_2(a_0)]$, and continuity gives $H_{a_0}(y)\leq0$.  The endpoints
cannot occur: $H_a(1)>0$ by the alternating-series test, while at
$y=q_2(a)$ the first three terms sum to $1$ and the remaining alternating
tail is larger than $-q_2/q_3>-1$.  Hence $1<y<q_2(a_0)$, and the sign
criterion gives $F_{a_0}\in\LP$.

For the numerical enclosure, all arithmetic below is exact.  Put
\[
 a_- =\frac{19821}{5000}=3.9642,
 \qquad S_{5,a}(x)=\sum_{k=0}^5
 \frac{(-1)^kx^k}{\prod_{j=1}^k(a^j+1)}.
\]
Map each of the $128$ equal subintervals of
$[a_-+1,a_-^2+1]$ to $[0,1]$ and express $S_{5,a_-}$ in the degree-five
Bernstein basis.  The accompanying verifier checks all $128\cdot6=768$
rational coefficients; every one is positive, and the least is approximately
$4.6399567291958566\cdot10^{-6}$.  Thus $S_{5,a_-}>0$ throughout the
interval.  The remaining alternating tail begins with a positive term and
has strictly decreasing magnitudes, so $F_{a_-}(-x)>S_{5,a_-}(x)>0$ there.
The sign criterion implies $F_{a_-}\notin\LP$.

At the upper rational endpoint
\[
 a_+=\frac{39643}{10000}=3.9643,
 \qquad x_+=\frac23(a_+^2+1),
\]
the same verifier obtains
\[
 S_{6,a_+}(x_+)=-\frac{N_+}{D_+}
 \approx-8.43053729684106\cdot10^{-6}<0,
\]
where the verifier prints the positive integers $N_+$ and $D_+$ in full.
The tail after the sixth section begins negative and has decreasing
magnitudes, hence $F_{a_+}(-x_+)<S_{6,a_+}(x_+)<0$.  Therefore
$F_{a_+}\in\LP$.  Both endpoint inequalities persist in a neighborhood by
continuity, giving the strict enclosure claimed.
\end{proof}

\section{Correction of the published upper bound}

The paper's Theorem~2(ii) asserts membership for every $a\geq3.91719$.
That assertion is false: Corollary~\ref{cor:threshold} gives, for example,
$F_{3.9642}\notin\LP$.  The source of the discrepancy is visible in the
denominator-clearing step of Lemma~5.  At the chosen point
$x=2(a^2+1)/3$, the displayed expression in the $q_j$ is correct, but after
multiplication by the common positive denominator its fifth contribution
must be
\[
 -96(a^2+1)^4(a^6+1),
\]
not the printed $-96(a^2+1)^4(a^5+1)$.  The printed expression therefore
exceeds the correct one by
$96a^5(a-1)(a^2+1)^4$.  Direct evaluation at $a=3.91719$ gives
$S_{6,a}(2(a^2+1)/3)\approx0.0099216822>0$, rather than the required
nonpositive value.

This correction does not affect Theorem~1, the sign characterization used
above.  It removes only the claimed $3.91719$ upper bound; the monotonicity
proof and exact section certificates replace it with a genuine sharp-threshold
theorem and a four-decimal localization.

\section{Verification and novelty scope}

Run the exact, dependency-free certificate with
\begin{verbatim}
python code/exact_certificate.py
\end{verbatim}
It uses only integer arithmetic and \texttt{fractions.Fraction}; no
floating-point comparison enters either endpoint proof.  The central
monotonicity argument is analytic and does not depend on this computation.

A bounded search through 11 August 2026 covered the run's result and attempt
indexes and web searches for the exact title, arXiv identifier, conjecture,
parameter monotonicity, and the constants $3.91719$ and $3.9642$.  It found
the original paper and related later work, but no later source resolving this
conjecture or correcting Theorem~2(ii).  Novelty confidence is moderate
pending specialist review; mathematical confidence is high, with the main
dependency made explicit above.

\begin{thebibliography}{9}
\bibitem{Nguyen}
T. H. Nguyen,
\emph{On the conditions for a special entire function relative to the
partial theta-function and the Euler function to belong to the
Laguerre--P\'olya class}, arXiv:1912.10035 (2019).
\url{https://arxiv.org/abs/1912.10035}.
\end{thebibliography}

\end{document}
